\section*{Chapitre \ref{ch:s2:mortgage-and-structured-credit-strategies} --- Stratégies sur crédits hypothécaires et crédits structurés}

\begin{solution}{exo:s2:mortgage-and-structured-credit-strategies:1}
$26.05 \times 0.358 = 9.32$ par 100 de valeur nominale.
\end{solution}

\begin{solution}{exo:s2:mortgage-and-structured-credit-strategies:2}
Dans la vue, $148 - 59 = 88$ points de base (à partir des valeurs non arrondies)~; sans essoufflement, $-326 - 42 = -368$ points de base. Les mêmes prix donnent un
gain ou une perte importants selon le modèle.
\end{solution}

\begin{solution}{exo:s2:mortgage-and-structured-credit-strategies:3}
La tranche représente 3~\% de l'indice, soit 30 millions~; $30 \times 0.1169 = 3.51$ millions.
\end{solution}

\begin{solution}{exo:s2:mortgage-and-structured-credit-strategies:4}
Le strip intérêts seuls (IO) ne perçoit des intérêts que sur le principal encore en circulation. Quand les taux baissent, les emprunteurs se refinancent, le principal disparaît et les intérêts
avec lui, si bien que l'IO perd de la valeur. Sa valeur augmente donc avec les taux, comme une position courte en obligations, et couvrir une exposition courte sur obligations revient à acheter des obligations.
\end{solution}

\begin{solution}{exo:s2:mortgage-and-structured-credit-strategies:5}
Gabaix, Krishnamurthy et Vigneron ont établi que le titre est évalué comme si l'investisseur marginal était un arbitragiste spécialisé dans les MBS, pour qui le risque de remboursement anticipé ne
se diversifie pas~: ce qui est neutre pour l'économie se concentre dans les bilans de quelques acteurs qui détiennent et couvrent les crédits hypothécaires.
\end{solution}

\begin{solution}{exo:s2:mortgage-and-structured-credit-strategies:6}
Une tranche senior ne subit de perte que lorsque de nombreux noms font défaut ensemble, si bien que sa perte attendue dépend de la corrélation, difficile à estimer, et non de la
perte attendue du portefeuille à elle seule. Dans le tableau du chapitre, la perte attendue de la tranche super senior est dix-huit fois plus élevée pour une corrélation de 0.45 que
de 0.15 pour le même portefeuille~; Coval, Jurek et Stafford ont désigné cette fragilité des notations face à de faibles erreurs comme l'une des causes de l'effondrement de la
finance structurée.
\end{solution}

\begin{solution}{exo:s2:mortgage-and-structured-credit-strategies:7}
Le prix de l'IO passe de 26.05 à 26.85, le gain du pari passe de 35.8~\% à 33.5~\% et la perte sans essoufflement passe de 25.5~\% à 26.5~\%. Une
volatilité plus forte envoie davantage de trajectoires vers un refinancement massif, où l'hypothèse d'essoufflement erronée coûte le plus, si bien que l'avantage du pari se réduit face à sa perte potentielle~;
un pari sur les remboursements anticipés est aussi un pari sur la volatilité des taux.
\end{solution}

\begin{solution}{exo:s2:mortgage-and-structured-credit-strategies:8}
L'OAS est un résultat du modèle~: aux mêmes prix, un modèle sans essoufflement attribue au coupon de 7.5~\% un OAS de $-326$ points de base et au coupon de 4.0~\% un OAS de
42, soit un écart de $-368$. Rien n'est verrouillé~; les 88 points de base ne sont gagnés que si les emprunteurs se comportent comme le prévoit le modèle du pari, et les
ratios de couverture dépendent eux aussi du modèle.
\end{solution}

\begin{solution}{pb:s2:mortgage-and-structured-credit-strategies:1}
\begin{enumerate}
\item Une position en titres adossés à des créances hypothécaires dont la valeur dépend des vitesses de remboursement anticipé, prise parce que le modèle de remboursement anticipé du trader diffère de celui du
marché, avec le risque de taux couvert.
\item Rotation du parc selon la rampe PSA, une courbe en S de refinancement avec un CPR maximal de 45~\%, essoufflement (burnout) $e^{-0.25c}$~; un portefeuille au taux de 6.5~\% qui sert 6.0~\%, avec 29
années restantes, sur des trajectoires de Hull-White autour d'une courbe plate à 4~\%, à un OAS de 50 points de base, alors que les nouveaux crédits coûtent environ 5.8~\%.
\item Le recul du refinancement chez les emprunteurs ayant déjà laissé passer des incitations~; le pari suppose un essoufflement de 1.0 au lieu de 0.25.
\item Le risque de remboursement anticipé, neutre dans l'ensemble, est évalué, mieux expliqué par un risque commun à l'ensemble du marché des MBS, comme si l'investisseur marginal était un
arbitragiste spécialisé.
\item 26.05 par 100 de valeur nominale~; couvert par l'achat de 0.51 Treasury à dix ans de valeur nominale 100.
\item Taux inchangés~: 0.0~\% (marché), $35.8~\%$ (pari), $16.4~\%$ (refinancement plus lent), $-25.5~\%$ (sans essoufflement)~; taux inférieurs de 100 points de base~:
10.1~\%, 63.6~\%, 31.1~\% et $-36.0~\%$.
\item La baisse des taux accroît l'incitation, et les modèles divergent sur ce que font les emprunteurs~; la hausse la supprime, et tous les modèles prédisent la
même rotation.
\item Aux prix du marché, son OAS passe de 59 points de base pour le coupon de 4.0~\% à 148 pour celui de 7.5~\%~; sans essoufflement, il passe de 42 à
$-326$.
\item Négocier des tranches les unes contre les autres, contre l'indice ou contre les noms, parce que son modèle de corrélation ou de pertes du portefeuille diffère de celui des prix.
\item La tranche equity perd moins et les tranches senior perdent davantage à mesure que la corrélation augmente~; la perte attendue du portefeuille, 3.6~\%, ne change pas.
\item Les notations sont extrêmement fragiles face à de faibles erreurs sur les risques sous-jacents, et exposées au risque systématique~; de nombreux produits structurés se sont comportés comme
des obligations catastrophe économiques, pour une rémunération inférieure à celle d'alternatives comparables.
\item Pour une corrélation vraie de 0.15~: gain moyen de 3.1~\% de la tranche, perte dans 11.3~\% des portefeuilles, tranche anéantie dans 5.4~\%~; pour 0.30~: $-0.4~\%$, 14.5~\%,
9.4~\%~; pour 0.45~: $-1.4~\%$, 15.0~\%, 10.9~\%.
\item Couvert, l'IO gagne 35.8~\% de son prix si le pari sur l'essoufflement est juste, et perd 25.5~\% si les emprunteurs ne s'essoufflent jamais (63.6~\% et $-36.0~\%$ après
une baisse des taux de 100 points de base).
\item Le fait que le marché valorise le comportement des emprunteurs comme un risque, et non comme une fonction connue des taux~: un modèle unique à un seul OAS donnerait une droite plate.
\item La couverture est dimensionnée pour de petits mouvements~: quand les taux baissent fortement, la valeur de l'IO s'aplatit, car il reste peu à rembourser, tandis que le Treasury continue
à monter~; quand les taux montent fortement, l'IO continue de gagner tandis que les pertes du Treasury ralentissent.
\item Estimer le modèle sur des cohortes et des périodes non utilisées pour l'ajuster, avec les données disponibles à l'époque, et comparer les vitesses prédites et réalisées par
coupon, ancienneté et taille du prêt.
\item Dans le cas où le modèle se trompe~: la perte sans essoufflement après un rallye obligataire, et non le gain espéré.
\item La corrélation de défaut n'apparaît que dans de rares vagues de défauts systématiques~; quelques années sans une telle vague ne permettent pas de distinguer 0.15 de 0.30.
\item La valeur relative des OAS entre coupons~: un classement transversal selon une caractéristique, la duration étant couverte.
\item L'option de remboursement de l'emprunteur~; l'avantage du trader est de savoir à quel point elle sera mal exercée.
\end{enumerate}
\end{solution}

\begin{solution}{iq:s2:mortgage-and-structured-credit-strategies:1}
Quand les taux baissent, les remboursements anticipés s'accélèrent et le prix du titre monte moins que celui d'une obligation~; quand les taux montent, les remboursements ralentissent et sa duration s'allonge,
si bien qu'il baisse davantage~: son prix est une fonction concave des taux, car le détenteur est vendeur de l'option de refinancement des emprunteurs.
\end{solution}

\begin{solution}{iq:s2:mortgage-and-structured-credit-strategies:2}
Séparer la rotation du parc (ancienneté, saison de l'année, prix immobiliers) du refinancement (incitation, taille du prêt, score de crédit, ratio prêt-valeur, essoufflement, effets
médiatiques), ainsi que les remboursements partiels et les défauts~; ajuster sur des données par prêt et par cohorte, et tester hors échantillon (out-of-sample) à travers les vagues de refinancement.
\end{solution}

\begin{solution}{iq:s2:mortgage-and-structured-credit-strategies:3}
Les remboursements anticipés s'accélèrent et l'IO perd de la valeur~; la couverture par le Treasury gagne. L'ampleur dépend des vitesses, il faut donc vérifier les vitesses réalisées par rapport au modèle et
rééquilibrer la couverture à mesure que la duration de l'IO change, en la raccourcissant à mesure que l'IO perd de la valeur.
\end{solution}

\begin{solution}{iq:s2:mortgage-and-structured-credit-strategies:4}
Sensibilité à la corrélation, saut au défaut de noms individuels, scénarios sur le niveau des spreads pour l'ensemble de l'indice, hypothèses de recouvrement, et la queue des pertes de portefeuille
issue d'une simulation, au-delà des deltas du premier ordre.
\end{solution}

\begin{solution}{iq:s2:mortgage-and-structured-credit-strategies:5}
Partager les trajectoires de taux entre portefeuilles (nombres aléatoires communs, générés une seule fois), vectoriser les flux de trésorerie entre portefeuilles et trajectoires, utiliser moins de trajectoires avec
réduction de variance, mettre en cache les termes du modèle de remboursement anticipé, paralléliser sur les cœurs, et ne recalculer que les portefeuilles dont les entrées ont changé.
\end{solution}

\begin{solution}{iq:s2:mortgage-and-structured-credit-strategies:6}
La perte du portefeuille sachant le facteur $Z$ est $(1-R)\,\Phi\bigl((\Phi^{-1}(p) - \sqrt\rho Z)/\sqrt{1-\rho}\bigr)$~; son espérance sur $Z$ est
$(1-R)\,P(\sqrt\rho Z + \sqrt{1-\rho}\,\varepsilon < \Phi^{-1}(p)) = (1-R)p$, quel que soit $\rho$. La perte d'une tranche est une fonction non linéaire et plafonnée de la
perte du portefeuille, si bien que son espérance dépend de toute la distribution de cette perte, que la corrélation étale.
\end{solution}
